
\documentclass{article}
\usepackage{colortbl}
\usepackage{makecell}
\usepackage{multirow}
\usepackage{supertabular}

\begin{document}

\newcounter{utterance}

\centering \large Interaction Transcript for game `cladder', experiment `full\_v1.5\_default', episode 9064 with Qwen3.5{-}9B.
\vspace{24pt}

{ \footnotesize  \setcounter{utterance}{1}
\setlength{\tabcolsep}{0pt}
\begin{supertabular}{c@{$\;$}|p{.15\linewidth}@{}p{.15\linewidth}p{.15\linewidth}p{.15\linewidth}p{.15\linewidth}p{.15\linewidth}}
    \# & \multicolumn{2}{c}{Player} && \multicolumn{2}{c}{Game Master} \\
    \hline

    \theutterance \stepcounter{utterance}  
    & & & \multicolumn{4}{p{0.6\linewidth}}{
        \cellcolor[rgb]{0.9,0.9,0.9}{
            \makecell[{{p{\linewidth}}}]{
                \texttt{\tiny{[P1$\langle$GM]}}
                \texttt{You are an expert in causal inference. The following question is not a typical commonsense query, but rather a meticulously designed question created by a professor specializing in causal inference, intended to assess the students' mastery of the course content.} \\
\texttt{Imagine a self{-}contained, hypothetical world with only the following conditions, and without any unmentioned factors or causal relationships: Pexu has a direct effect on hwax. Hwax has a direct effect on rukz. Method 1: We look at how pexu correlates with rukz case by case according to hwax. Method 2: We look directly at how pexu correlates with rukz in general. To understand how pexu affects rukz, is it more correct to use the Method 1 than Method 2?} \\
\texttt{Start your answer with 'Yes' or 'No', followed by additional reasoning or evidence to support your explanation.} \\
            }
        }
    }
     \\ \\

    \theutterance \stepcounter{utterance}  
    & \multicolumn{4}{p{0.6\linewidth}}{
        \cellcolor[rgb]{0.561,0.737,0.561}{
            \makecell[{{p{\linewidth}}}]{
                \texttt{\tiny{[P1$\rangle$GM]}}
                \texttt{No} \\
\\ 
\texttt{In this hypothetical world, the causal structure forms a simple mediation chain: \$Pexu \textbackslash{}rightarrow Hwax \textbackslash{}rightarrow Rukz\$. To understand the effect of \$Pexu\$ on \$Rukz\$, we must distinguish between the **total effect** and the **direct effect**.} \\
\\ 
\texttt{Method 2, which looks directly at the correlation between \$Pexu\$ and \$Rukz\$ in general, estimates the **total effect**. This captures the full impact of \$Pexu\$ on \$Rukz\$ by summing both the direct path (if any existed) and the indirect path through \$Hwax\$. In this specific scenario where \$Pexu\$ only affects \$Rukz\$ via \$Hwax\$, Method 2 correctly identifies the total causal influence.} \\
\\ 
\texttt{Method 1, described as looking at how \$Pexu\$ correlates with \$Rukz\$ "case by case according to \$Hwax\$," implies conditioning on or stratifying by the mediator \$Hwax\$. In causal inference, conditioning on a mediator (a variable on the causal pathway) blocks the path through which the treatment affects the outcome. This results in estimating the **direct effect** (which is zero in this specific chain, as there is no direct arrow from \$Pexu\$ to \$Rukz\$) or, if the conditioning is imperfect, a biased estimate of the total effect due to collider bias or simply isolating the wrong component.} \\
\\ 
\texttt{Therefore, to understand "how \$Pexu\$ affects \$Rukz\$" in the broad sense of the system's behavior, Method 2 is the correct approach to measure the total effect. Method 1 is incorrect for this goal because it artificially blocks the very mechanism (the effect through \$Hwax\$) that constitutes the relationship in this world. Using Method 1 would lead one to conclude that \$Pexu\$ has no effect on \$Rukz\$, which contradicts the premise that \$Pexu\$ has a direct effect on \$Hwax\$ and \$Hwax\$ has a direct effect on \$Rukz\$.} \\
            }
        }
    }
    & & \\ \\

    \theutterance \stepcounter{utterance}  
    & & & \multicolumn{2}{p{0.3\linewidth}}{
        \cellcolor[rgb]{0.9,0.9,0.9}{
            \makecell[{{p{\linewidth}}}]{
                \texttt{\tiny{[GM$|$GM]}}
                \texttt{No} \\
            }
        }
    }
    & & \\ \\

    \theutterance \stepcounter{utterance}  
    & & & \multicolumn{2}{p{0.3\linewidth}}{
        \cellcolor[rgb]{0.9,0.9,0.9}{
            \makecell[{{p{\linewidth}}}]{
                \texttt{\tiny{[GM$|$GM]}}
                \texttt{no} \\
            }
        }
    }
    & & \\ \\

    \theutterance \stepcounter{utterance}  
    & & & \multicolumn{2}{p{0.3\linewidth}}{
        \cellcolor[rgb]{0.9,0.9,0.9}{
            \makecell[{{p{\linewidth}}}]{
                \texttt{\tiny{[GM$|$GM]}}
                \texttt{game\_result = WIN} \\
            }
        }
    }
    & & \\ \\

\end{supertabular}
}

\end{document}
